\documentclass[12pt]{article}
\usepackage[margin=1in]{geometry}
\usepackage{amsmath}
\usepackage{fancyhdr}
\fancypagestyle{plain}{%
  \fancyhf{}
  \fancyhead[L]{POLISCI 450A}
  \fancyhead[R]{Fall 2026}
  %\fancyfoot[C]{\thepage}
}

\usepackage{titlesec}
\titleformat{\section}{\large\bfseries}{\thesection}{1em}{}
\usepackage{titling}
\setlength{\droptitle}{-3em}   % move title up


\title{\Large Lecture 3: What Is Regression?}
\author{}
\date{}
\begin{document}
\maketitle
\vspace{-6em}

\noindent \textbf{Goal:} Introduce regression as estimating $E[Y|X]$; compare nonparametric and parametric approaches.\medskip

\begin{itemize}\setlength{\itemsep}{1em}
\item What is regression?
\begin{itemize}
\item Conditioning on predictors
\item Modeling conditional mean $E[Y|X]$
\item Prediction versus causal uses
\item Estimator: $\widehat{E}[Y|X]$
\end{itemize}
\item Nonparametric regression: discrete $X$
\begin{itemize}
\item Group means by education levels
\item Estimate $E[Y|X=x]$ within categories
\end{itemize}
\item Nonparametric regression: continuous $X$
\begin{itemize}
\item Kernel regression idea
\item Uniform kernel local averages
\item Bandwidth controls smoothing
\item Weighted kernels (Epanechnikov \& Gaussian)
\item Bias–variance tradeoff
\end{itemize}
\item Linear regression
\begin{itemize}
\item Parametric form $E[Y|X]=\beta_0+\beta_1 X$
\item Interpret intercept and slope
\item Justifications and limitations
\end{itemize}
\item Local(ly) linear regressions
\begin{itemize}
\item LOWESS/LOESS local linear smoothing
\item Bandwidth and weights in practice
\item GAM as flexible alternative
\end{itemize}
\end{itemize}


% End outline
\clearpage


%\section*{Worksheet}


% \begin{enumerate}\itemsep1em
% \item Which of $\hat\mu_{1}, \hat\mu_{2}, \hat\mu_{3}, \hat\mu_{4}$ are unbiased? Show your work.
% \vspace{20em}

% \item Derive the variances of $\hat\mu_{1}, \hat\mu_{2}, \hat\mu_{3}, \hat\mu_{4}$. Show your work.
% \vspace{20em}
% \clearpage 

% \item Prove: $MSE(\hat{\theta}) = Var(\hat{\theta}) + \text{Bias}(\hat{\theta})^2.$
% \vspace{25em}


% \item Derive $MSE(\hat\mu_{1}), MSE(\hat\mu_{2}), MSE(\hat\mu_{3}), MSE(\hat\mu_{4})$
% \vspace{20em}


% \end{enumerate}


\end{document}
